\section{Построение конечного автомата по регулярному выражению}
\label{sec:RegexToFA}
\tikzsetfigurename{RegLang_RegexToFA_}

Конечные автоматы и регулярные выражения~--- два способа задать один и тот же класс языков.
Чтобы убедиться в их равной выразительной силе, начнём с построения конечного автомата по регулярному выражению.
Использование производных позволит нам сразу получить \emph{полный детерминированный автомат}%
\sidenote{Существуют и другие алгоритмы построения автомата по регулярному выражению, например, алгоритм Томпсона~\cite{10.1145/363347.363387} или алгоритм Глушкова~\cite{Glushkov1961}.
Каждый из них имеет свои особенности: некоторые строят недетерминированные автоматы, другие используют $\varepsilon$-переходы.}.

Состояниям автомата будем сопоставлять регулярные выражения.
Если выбрать некоторое состояние автомата в качестве стартового, то получившийся автомат будет распознавать язык, задаваемый соответствующим регулярным выражением.

Чтобы избежать дублирования состояний, необходимо проверять эквивалентность языков, которые задаются этими выражениями.
Хотя задача проверки эквивалентности регулярных языков разрешима, на практике в данном алгоритме используют более простую \emph{структурную эквивалентность} регулярных выражений%
\sidenote{При проверке структурной эквивалентности обычно учитывают ассоциативность, коммутативность и идемпотентность некоторых операций, чтобы избежать избыточных различий.}~\sidecite{OWENS_REPPY_TURON_2009}.


Таким образом, можно сформулировать следующий алгоритм построения автомата по регулярному выражению $R$.

\begin{enumerate}
    \item В качестве стартового состояния берём выражение $R$ и добавляем его в рабочее множество $W$.
    \item Пока $W$ не пусто:
    \begin{enumerate}
        \item Извлекаем очередное состояние $q$, которому соответствует регулярное выражение $R_q$.
        \item Для каждого символа $a \in \Sigma$ вычисляем производную $\partial_a(R_q)$.
        \item Если полученное выражение задаёт новый (ранее не встречавшийся) язык, добавляем соответствующее состояние.
        \item Если $\emph{IsNull}(R_q) = \emph{true}$, помечаем состояние как финальное.
        \item Добавляем ребро $q \xrightarrow{a} q'$ для каждого вычисленного перехода.
    \end{enumerate}
\end{enumerate}

В результате получим полный детерминированный конечный автомат\sidenote{Формально, если честно проверять равенство языков, можно получить минимальный, но на практике это сложно. Но часто всё же достаточно близко к минимальному получается.}.

\begin{example}
    Рассмотрим построение автомата для регулярного выражения $a(a \mid (b \ b))^*$.

    Стартовому состоянию соответствует выражение $a(a \mid (b \ b))^*$.
    Так как $\emph{IsNull}(a(a \mid (b \ b))^*) = \emph{false}$, это состояние не является финальным.
    Вычислим производные по всем символам из алфавита:
    \begin{align*}
        \partial_a(a(a \mid (b \ b))^*) & = (a \mid (b \ b))^*, \\
        \partial_b(a(a \mid (b \ b))^*) & = \varnothing.
    \end{align*}
    В результате появляются два новых состояния: $(a \mid (b \ b))^*$ и $\varnothing$.
    Второе состояние соответствует пустому языку и является \emph{поглощающим} (или \emph{дьявольским}) состоянием, а первое~--- финальным, так как $\emph{IsNull}((a \mid (b \ b))^*) = \emph{true}$.
    На рисунке~\ref{fig:regexp_to_dfa_example_step_1} показан автомат после первого шага алгоритма.

    \begin{marginfigure}
        \begin{center}
            \scalebox{0.64}{
                        \begin{tikzpicture}
                \node[elliptic state,initial] (q_0) {$a(a \mid (b \ b))^*$};
                \node[elliptic state,accepting] (q_1) [right = of q_0] {$(a \mid (b \ b))^*$};
                \node[state] (q_2) [above = of q_0] {$\varnothing$};
                \path[->]
                    (q_0) edge[bend right, below]   node {$a$} (q_1)
                    (q_0) edge[bend left, left]   node {$b$} (q_2);
            \end{tikzpicture}

            }
        \end{center}
        \caption{Построение автомата по регулярному выражению с помощью производных: первый шаг}
        \label{fig:regexp_to_dfa_example_step_1}
    \end{marginfigure}

    На следующем шаге вычислим производные для состояния $(a \mid (b \ b))^*$:
    \begin{align*}
        \partial_a((a \mid (b \ b))^*) & = (a \mid (b \ b))^*, \\
        \partial_b((a \mid (b \ b))^*) & = b(a \mid (b \ b))^*.
    \end{align*}
    Таким образом, появляется новое состояние $b(a \mid (b \ b))^*$, а автомат принимает вид, показанный на рисунке~\ref{fig:regexp_to_dfa_example_step_2}.

    \begin{marginfigure}
        \begin{center}
            \scalebox{0.64}{
            \input{part_02_Foundations/chapter_05_RegularLanguages/figures/04_RegexToFA/fig_regexp_to_dfa_example_step_2.tex}
            }
        \end{center}
        \caption{Построение автомата по регулярному выражению с помощью производных: второй шаг}
        \label{fig:regexp_to_dfa_example_step_2}
    \end{marginfigure}

    Осталось обработать состояние $b(a \mid (b \ b))^*$:
    \begin{align*}
        \partial_a(b(a \mid (b \ b))^*) & = \varnothing, \\
        \partial_b(b(a \mid (b \ b))^*) & = (a \mid (b \ b))^*.
    \end{align*}
    В результате добавляются новые переходы, и автомат принимает окончательный вид (рисунок~\ref{fig:regexp_to_dfa_example_step_3}).

    \begin{marginfigure}
        \begin{center}
            \scalebox{0.64}{
                        \begin{tikzpicture}
                \node[elliptic state,initial] (q_0) {$a(a \mid (b \ b))^*$};
                \node[elliptic state,accepting] (q_1) [right = of q_0] {$(a \mid (b \ b))^*$};
                \node[state] (q_2) [above = of q_0] {$\varnothing$};
                \node[elliptic state] (q_3) [above = of q_1] {$b(a \mid (b \ b))^*$};
                \path[->]
                    (q_0) edge[bend right, below]   node {$a$} (q_1)
                    (q_0) edge[bend left, left]   node {$b$} (q_2)
                    (q_1) edge[bend left, left]   node {$b$} (q_3)
                    (q_2) edge[loop left, left]   node {$b,a$} (q_2)
                    (q_1) edge[loop below, above]   node {$a$} (q_1)
                    (q_3) edge[bend left, right]   node {$b$} (q_1)
                    (q_3) edge[bend left, below]   node {$a$} (q_2);
            \end{tikzpicture}

            }
        \end{center}
        \caption{Построение автомата по регулярному выражению с помощью производных: результирующий автомат}
        \label{fig:regexp_to_dfa_example_step_3}
    \end{marginfigure}
\end{example}

В результате построения получен автомат, близкий по структуре к автомату из примера~\ref{exmpl:DFA}.
Разница лишь в том, что приведённый здесь автомат является полным.
Само же сходство не случайно: автоматы действительно задают один и тот же регулярный язык.
